\documentclass[12pt,a4paper]{article}
\usepackage[utf8]{inputenc}
\usepackage[greek,english]{babel}
\usepackage{geometry}
\geometry{a4paper, margin=1in}
\usepackage{amsmath}
\usepackage{amssymb}
\usepackage{graphicx}
\usepackage{listings}
\usepackage{xcolor}
\usepackage{hyperref}
\usepackage{fancyhdr}
\usepackage{float}
\usepackage{caption}

% Listings configuration for Python code
\definecolor{codegreen}{rgb}{0,0.6,0}
\definecolor{codegray}{rgb}{0.5,0.5,0.5}
\definecolor{codepurple}{rgb}{0.58,0,0.82}
\definecolor{backcolour}{rgb}{0.95,0.95,0.92}

\lstdefinestyle{pythonstyle}{
    backgroundcolor=\color{backcolour},   
    commentstyle=\color{codegreen},
    keywordstyle=\color{magenta},
    numberstyle=\tiny\color{codegray},
    stringstyle=\color{codepurple},
    basicstyle=\ttfamily\footnotesize,
    breakatwhitespace=false,         
    breaklines=true,                 
    captionpos=b,                    
    keepspaces=true,                 
    numbers=left,                    
    numbersep=5pt,                  
    showspaces=false,                
    showstringspaces=false,
    showtabs=false,                  
    tabsize=2,
    language=Python,
    frame=single
}

\lstset{style=pythonstyle}

% Page style
\pagestyle{fancy}
\fancyhf{}
\rhead{\selectlanguage{greek}Τεκμηρίωση Κώδικα Ανάλυσης FEM}
\lhead{\selectlanguage{greek}Ρακοβάλης Παύλος - 6931}
\rfoot{\selectlanguage{greek}Σελίδα \thepage}

\title{\selectlanguage{greek}\textbf{Μέθοδος Πεπερασμένων Στοιχείων (FEM)\\
Επιλύτης 3D Στοιχείων Δοκού}\\
\large Τεκμηρίωση Κώδικα και Αναφορά Υλοποίησης}
\author{\selectlanguage{greek}Ρακοβάλης Παύλος\\
Αριθμός Μητρώου: 6931\\
\small Μάθημα: Μέθοδος Πεπερασμένων Στοιχείων}
\date{\selectlanguage{greek}Δεκέμβριος 2025}

\begin{document}
\selectlanguage{greek}

\maketitle
\thispagestyle{empty}

\newpage
\tableofcontents
\newpage

\section{Εισαγωγή}

Η παρούσα αναφορά τεκμηριώνει μια ολοκληρωμένη υλοποίηση λογισμικού Μεθόδου Πεπερασμένων Στοιχείων (FEM) για την ανάλυση τρισδιάστατων κατασκευών δοκών. Το λογισμικό ακολουθεί μια αρθρωτή αρχιτεκτονική με τρία κύρια συστατικά:

\begin{itemize}
    \item \textbf{Προεπεξεργαστής} (\texttt{preprocessor.py}): Δημιουργία γεωμετρίας, οριακές συνθήκες και εξαγωγή δεδομένων
    \item \textbf{Επιλύτης} (\texttt{solver.py}): Συναρμολόγηση και επίλυση συστήματος FEM
    \item \textbf{Μετεπεξεργαστής} (\texttt{postprocessor.py}): Οπτικοποίηση και ανάλυση αποτελεσμάτων
\end{itemize}

Η υλοποίηση χειρίζεται τρισδιάστατα στοιχεία δοκών με 6 βαθμούς ελευθερίας (DOF) ανά κόμβο: τρεις μεταφορικές μετατοπίσεις $(U_x, U_y, U_z)$ και τρεις περιστροφικές μετατοπίσεις $(R_x, R_y, R_z)$.

\subsection{Κύρια Χαρακτηριστικά}

\begin{itemize}
    \item Υποστήριξη για σύνθετες τρισδιάστατες κατασκευές πλαισίων και δοκών
    \item Αυτόματη δημιουργία στοιχείων με παράλληλα προς άξονες και διαγώνια ενισχυτικά στοιχεία
    \item Γεωμετρικοί μετασχηματισμοί (περιστροφή γύρω από αυθαίρετους άξονες)
    \item Ιδιότητες δοκών κυκλικής διατομής
    \item Ολοκληρωμένος χειρισμός οριακών συνθηκών
    \item Διαδραστική οπτικοποίηση με χρήση Plotly
    \item Αρθρωτή, επεκτάσιμη αρχιτεκτονική
\end{itemize}

\section{Μαθηματική Θεμελίωση}

\subsection{Διατύπωση Τρισδιάστατου Στοιχείου Δοκού}

Για ένα τρισδιάστατο στοιχείο δοκού με 6 DOF ανά κόμβο, ο τοπικός πίνακας δυσκαμψίας είναι $12 \times 12$. Το στοιχείο περιλαμβάνει:

\begin{itemize}
    \item \textbf{Αξονική παραμόρφωση}: Σχετίζεται με το μέτρο ελαστικότητας $E$ και την εγκάρσια διατομή $A$
    \item \textbf{Στρέψη}: Σχετίζεται με το μέτρο διάτμησης $G$ και τη σταθερά στρέψης $J$
    \item \textbf{Κάμψη}: Σχετίζεται με το $E$ και τις ροπές αδράνειας $I_y$, $I_z$
\end{itemize}

Οι συντελεστές δυσκαμψίας υπολογίζονται ως:
\begin{align}
k_{αξονική} &= \frac{EA}{L} \\
k_{στρέψη} &= \frac{GJ}{L} \\
k_{κάμψη,y} &= \frac{12EI_y}{L^3}, \quad \frac{6EI_y}{L^2}, \quad \frac{4EI_y}{L}, \quad \frac{2EI_y}{L} \\
k_{κάμψη,z} &= \frac{12EI_z}{L^3}, \quad \frac{6EI_z}{L^2}, \quad \frac{4EI_z}{L}, \quad \frac{2EI_z}{L}
\end{align}

\subsection{Μετασχηματισμός Συντεταγμένων}

Ο μετασχηματισμός από τοπικές σε ολικές συντεταγμένες χρησιμοποιεί συνημίτονα κατεύθυνσης:
\begin{equation}
\mathbf{K}_{ολικός} = \mathbf{T}^T \mathbf{K}_{τοπικός} \mathbf{T}
\end{equation}

όπου $\mathbf{T}$ είναι ο πίνακας μετασχηματισμού $12 \times 12$ που κατασκευάζεται από τον πίνακα περιστροφής $3 \times 3$.

\section{Αρχιτεκτονική Λογισμικού}

\subsection{Επισκόπηση Μονάδων}

Το λογισμικό είναι οργανωμένο σε τέσσερις κύριες μονάδες Python:

\begin{enumerate}
    \item \textbf{FEM\_Main.py}: Κύριο σενάριο εκτέλεσης
    \item \textbf{preprocessor.py}: Δημιουργία γεωμετρίας και προετοιμασία δεδομένων
    \item \textbf{solver.py}: Υλοποίηση επιλύτη FEM
    \item \textbf{postprocessor.py}: Οπτικοποίηση αποτελεσμάτων
    \item \textbf{geometry\_utils.py}: Βοηθητικά προγράμματα γεωμετρικού μετασχηματισμού
\end{enumerate}

\begin{figure}[H]
\centering
\begin{verbatim}
    +------------------+
    |   FEM_Main.py    |
    +------------------+
             |
             v
    +------------------+       +-------------------+
    | preprocessor.py  | ----> | structure.dat     |
    +------------------+       +-------------------+
             |                          |
             v                          v
    +------------------+       +-------------------+
    |   solver.py      | ----> | structure.res     |
    +------------------+       +-------------------+
             |                          |
             v                          v
    +------------------+       +-------------------+
    | postprocessor.py | ----> | Διαδραστικά Plots |
    +------------------+       +-------------------+
\end{verbatim}
\caption{Ροή εργασίας λογισμικού και ροή δεδομένων}
\end{figure}

\section{Μονάδα Προεπεξεργαστή}

Ο προεπεξεργαστής (\texttt{preprocessor.py}) είναι υπεύθυνος για τη δημιουργία της δομικής γεωμετρίας, τον ορισμό οριακών συνθηκών και την εξαγωγή δεδομένων σε τυποποιημένη μορφή.

\subsection{Δομή Κλάσης}

Η κύρια κλάση είναι η \texttt{FEMPreProcessor}:

\begin{lstlisting}[caption={Αρχικοποίηση Κλάσης FEMPreProcessor}]
class FEMPreProcessor:
    """Handles geometry creation and boundary conditions for FEM analysis"""
    
    def __init__(self):
        self.points = pd.DataFrame(columns=['Node Number', 'X', 'Y', 'Z'])
        self.elements = pd.DataFrame(columns=['Element Number', 'Node1', 'Node2', 'E', 'V'])
        self.boundary_conditions = {}  # node_id: {'Ux': 0/1, 'Uy': 0/1, 'Uz': 0/1, 'Rx': 0/1, 'Ry': 0/1, 'Rz': 0/1}
        self.loads = {}  # node_id: {'Fx': value, 'Fy': value, 'Fz': value, 'Mx': value, 'My': value, 'Mz': value}
        self.materials = {'STEEL': {'E': 210e9, 'nu': 0.3}}
        
        # Beam properties (circular cross-section)
        self.diameter = 0.020  # 20mm diameter circular cross-section
\end{lstlisting}

\subsection{Ιδιότητες Κυκλικής Διατομής}

Για κυκλικές διατομές, οι γεωμετρικές ιδιότητες υπολογίζονται ως εξής:

\begin{lstlisting}[caption={Υπολογισμός Ιδιοτήτων Κυκλικής Διατομής}]
def _calculate_circular_properties(self, diameter):
    """Calculate geometric properties for circular cross-section"""
    r = diameter / 2.0
    
    # Cross-sectional area
    A = np.pi * r**2
    
    # Second moment of area (Iz = Iy for circular section)
    I = (np.pi * r**4) / 4.0
    Iz = I
    Iy = I
    
    # Torsional constant (polar moment of inertia)
    J = (np.pi * r**4) / 2.0
    
    return {'A': A, 'Iz': Iz, 'Iy': Iy, 'J': J, 'diameter': diameter, 'radius': r}
\end{lstlisting}

Αυτοί οι τύποι υλοποιούν τις τυπικές σχέσεις για κυκλικές διατομές:
\begin{align}
A &= \pi r^2 \\
I_y = I_z &= \frac{\pi r^4}{4} \\
J &= \frac{\pi r^4}{2} = 2I
\end{align}

\subsection{Δημιουργία Γεωμετρίας}

Η δημιουργία γεωμετρίας περιλαμβάνει πολλαπλά βήματα, συμπεριλαμβανομένου του ορισμού κόμβων, της δημιουργίας στοιχείων και των γεωμετρικών μετασχηματισμών.

\begin{lstlisting}[caption={Ορισμός Παραμέτρων Γεωμετρίας}]
def _define_geometry_parameters(self):
    """Define geometry-specific parameters (not saved to .dat file)"""
    # COMPLEX GEOMETRY: Original crane/truss structure
    # Geometry parameters for this specific crane/truss structure
    self.L = 1.5 * 1.13      # Element length (m)
    self.A = 1.2 * 1.69      # Base dimension (m)
    self.B = 1.93            # Additional parameter B (m)
    self.phi = 60 + 9.1      # Rotation angle (degrees)
    self.A_0 = 6 * (0.5 + 0.13)  # Base cross-section parameter (dimensionless)
\end{lstlisting}

\subsection{Δημιουργία Στοιχείων}

Τα στοιχεία δημιουργούνται αυτόματα με βάση μοτίβα συνδεσιμότητας:

\begin{lstlisting}[caption={Δημιουργία Στοιχείων Παράλληλων προς Άξονες}]
def _create_axis_aligned_elements(self, element_counter, tolerance, max_length):
    """Create elements parallel to X, Y, or Z axes"""
    axis_aligned_count = 0
    for i in range(len(self.points)):
        node1 = self.points.iloc[i]
        node1_num = int(node1['Node Number'])
        
        for j in range(i + 1, len(self.points)):
            node2 = self.points.iloc[j]
            node2_num = int(node2['Node Number'])
            
            dx = abs(node2['X'] - node1['X'])
            dy = abs(node2['Y'] - node1['Y'])
            dz = abs(node2['Z'] - node1['Z'])
            
            element_length = np.sqrt(dx**2 + dy**2 + dz**2)
            
            if element_length > max_length:
                continue
            
            # Check axis alignment
            is_x_aligned = dx > tolerance and dy < tolerance and dz < tolerance
            is_y_aligned = dx < tolerance and dy > tolerance and dz < tolerance
            is_z_aligned = dx < tolerance and dy < tolerance and dz > tolerance
\end{lstlisting}

\subsection{Οριακές Συνθήκες και Φορτία}

Ο προεπεξεργαστής παρέχει μεθόδους για τον ορισμό οριακών συνθηκών και εφαρμοζόμενων φορτίων:

\begin{lstlisting}[caption={Προσθήκη Οριακών Συνθηκών}]
def add_boundary_condition(self, node_id, ux=1, uy=1, uz=1, rx=1, ry=1, rz=1):
    """
    Add boundary condition for a node
    
    Parameters:
    -----------
    node_id : int
        Node number (1-based)
    ux, uy, uz : int
        Translation constraints (1=fixed, 0=free)
    rx, ry, rz : int
        Rotation constraints (1=fixed, 0=free)
    """
    self.boundary_conditions[node_id] = {
        'Ux': ux, 'Uy': uy, 'Uz': uz,
        'Rx': rx, 'Ry': ry, 'Rz': rz
    }
\end{lstlisting}

\begin{lstlisting}[caption={Προσθήκη Φορτίων}]
def add_load(self, node_id, fx=0.0, fy=0.0, fz=0.0, mx=0.0, my=0.0, mz=0.0):
    """
    Add load (forces and moments) to a node
    
    Parameters:
    -----------
    node_id : int
        Node number (1-based)
    fx, fy, fz : float
        Force components in Newtons
    mx, my, mz : float
        Moment components in Newton-meters
    """
    self.loads[node_id] = {
        'Fx': fx, 'Fy': fy, 'Fz': fz,
        'Mx': mx, 'My': my, 'Mz': mz
    }
\end{lstlisting}

\subsection{Εξαγωγή Δεδομένων}

Ο προεπεξεργαστής εξάγει όλα τα δομικά δεδομένα σε ένα αρχείο κειμένου σε δομημένη μορφή:

\begin{lstlisting}[caption={Εξαγωγή σε Αρχείο DAT}]
def export_to_file(self, filename='data/structure.dat'):
    """Export structure to text file for solver"""
    print("\nEXPORTING STRUCTURE DATA")
    print("="*80)
    
    # Ensure directory exists
    os.makedirs(os.path.dirname(filename), exist_ok=True)
    
    with open(filename, 'w') as f:
        # Write header
        f.write("# FEM Structure Data File\n")
        f.write("# Generated by FEM Pre-Processor (Beam Elements)\n")
        f.write("# Author: Rakovalis Pavlos 6931\n\n")
        
        # Write parameters
        f.write("PARAMETERS\n")
        props = self._calculate_circular_properties(self.diameter)
        f.write(f"{self.E} {props['A']} {props['Iz']} {props['Iy']} {props['J']} {self.G} {self.diameter}\n\n")
        
        # Write materials
        f.write("MATERIALS\n")
        for mat_name, props in self.materials.items():
            f.write(f"{mat_name} {props['E']} {props['nu']}\n")
        f.write("\n")
        
        # Write nodes
        f.write("NODES\n")
        for idx in range(len(self.points)):
            node = self.points.iloc[idx]
            f.write(f"{int(node['Node Number'])} {node['X']:.10f} {node['Y']:.10f} {node['Z']:.10f}\n")
        f.write("\n")
        
        # Write elements
        f.write("ELEMENTS\n")
        for idx in range(len(self.elements)):
            elem = self.elements.iloc[idx]
            f.write(f"{int(elem['Element Number'])} {int(elem['Node1'])} {int(elem['Node2'])} ")
            f.write(f"{elem['Element Cross Section']:.10e} {elem['E']:.10e} {elem['V']:.10f}\n")
        f.write("\n")
\end{lstlisting}

\section{Μονάδα Βοηθητικών Προγραμμάτων Γεωμετρίας}

Η μονάδα \texttt{geometry\_utils.py} παρέχει συναρτήσεις γεωμετρικού μετασχηματισμού.

\subsection{Περιστροφή Γύρω από τον Άξονα Y}

\begin{lstlisting}[caption={Συνάρτηση Περιστροφής Σημείων}]
def rotate_points_around_y(df: pd.DataFrame, angle_deg: float, origin=(0.0, 0.0, 0.0)) -> pd.DataFrame:
    """
    Rotate all points around the global Y-axis by angle_deg degrees.

    Args:
        df: DataFrame with columns ['X','Y','Z'] (and optionally 'Node Number']).
        angle_deg: Rotation angle in degrees (right-hand rule about +Y).
        origin: Pivot point (x0, y0, z0) to rotate about. Defaults to the global origin.

    Returns:
        A new DataFrame with rotated coordinates.
    """
    theta = np.deg2rad(angle_deg)
    c, s = np.cos(theta), np.sin(theta)
    # Rotation about Y: Ry = [[c,0,s],[0,1,0],[-s,0,c]]
    R = np.array([[c, 0.0, s],
                  [0.0, 1.0, 0.0],
                  [-s, 0.0, c]])

    coords = df[['X', 'Y', 'Z']].to_numpy(dtype=float)
    pivot = np.array(origin, dtype=float)
    rotated = (coords - pivot) @ R.T + pivot

    out = df.copy()
    out[['X', 'Y', 'Z']] = rotated
    return out
\end{lstlisting}

Ο πίνακας περιστροφής γύρω από τον άξονα Y δίνεται από:
\begin{equation}
\mathbf{R}_y(\theta) = \begin{bmatrix}
\cos\theta & 0 & \sin\theta \\
0 & 1 & 0 \\
-\sin\theta & 0 & \cos\theta
\end{bmatrix}
\end{equation}

\section{Μονάδα Επιλύτη}

Ο επιλύτης (\texttt{solver.py}) χειρίζεται τη συναρμολόγηση και την επίλυση του συστήματος FEM.

\subsection{Δομή Κλάσης}

\begin{lstlisting}[caption={Αρχικοποίηση Κλάσης FEMBeamSolver}]
class FEMBeamSolver:
    """Handles FEM system assembly and solution for 3D beam elements"""
    
    def __init__(self):
        self.nodes = pd.DataFrame()
        self.elements = pd.DataFrame()
        self.boundary_conditions = {}
        self.loads = {}
        self.materials = {}
        
        # Beam parameters
        self.E = 0.0
        self.A = 0.0
        self.Iz = 0.0
        self.Iy = 0.0
        self.J = 0.0
        self.G = 0.0
        self.diameter = 0.0
        
        # Solution arrays (6 DOF per node)
        self.dof_per_node = 6
        self.K_global = None
        self.F_global = None
        self.u_global = None
        self.reactions = None
\end{lstlisting}

\subsection{Υπολογισμός Πίνακα Μετασχηματισμού}

Ο πίνακας μετασχηματισμού συσχετίζει τοπικά και ολικά συστήματα συντεταγμένων:

\begin{lstlisting}[caption={Υπολογισμός Πίνακα Μετασχηματισμού}]
def _compute_transformation_matrix(self, x1, y1, z1, x2, y2, z2):
    """Compute 3x3 transformation matrix for beam local coordinate system"""
    # Element length and direction cosines
    dx = x2 - x1
    dy = y2 - y1
    dz = z2 - z1
    L = np.sqrt(dx**2 + dy**2 + dz**2)
    
    # Local x-axis (along beam element)
    x_local = np.array([dx/L, dy/L, dz/L])
    
    # Define auxiliary vector for determining local y and z axes
    # Use global Y-axis as reference, unless element is parallel to Y
    if abs(x_local[1]) > 0.9:  # Nearly vertical element
        aux = np.array([1.0, 0.0, 0.0])  # Use global X as auxiliary
    else:
        aux = np.array([0.0, 1.0, 0.0])  # Use global Y as auxiliary
    
    # Local z-axis (perpendicular to x_local and aux)
    z_local = np.cross(x_local, aux)
    z_local = z_local / np.linalg.norm(z_local)
    
    # Local y-axis (perpendicular to x and z)
    y_local = np.cross(z_local, x_local)
    y_local = y_local / np.linalg.norm(y_local)
    
    # 3x3 transformation matrix [x_local, y_local, z_local] as rows
    T = np.vstack([x_local, y_local, z_local])
    
    return T, L
\end{lstlisting}

Ο πίνακας μετασχηματισμού επεκτείνεται από $3 \times 3$ σε $12 \times 12$ για 6 DOF ανά κόμβο:

\begin{lstlisting}[caption={Επέκταση Πίνακα Μετασχηματισμού}]
def _expand_transformation_matrix(self, T):
    """Expand 3x3 transformation to 12x12 for 6 DOF per node"""
    T_12x12 = np.zeros((12, 12))
    # Apply transformation to each 3x3 block (translations and rotations)
    for i in range(4):
        T_12x12[3*i:3*i+3, 3*i:3*i+3] = T
    return T_12x12
\end{lstlisting}

\subsection{Τοπικός Πίνακας Δυσκαμψίας}

Ο τοπικός πίνακας δυσκαμψίας ενσωματώνει αξονικές, στρεπτικές και καμπτικές συνεισφορές:

\begin{lstlisting}[caption={Συναρμολόγηση Τοπικού Πίνακα Δυσκαμψίας (Μερική)}]
def _create_local_stiffness_matrix(self, L):
    """Create 12x12 local stiffness matrix for 3D beam element"""
    k_local = np.zeros((12, 12))
    
    # Axial stiffness (DOF 0 and 6)
    EA_L = self.E * self.A / L
    k_local[0, 0] = EA_L
    k_local[0, 6] = -EA_L
    k_local[6, 0] = -EA_L
    k_local[6, 6] = EA_L
    
    # Torsional stiffness (DOF 3 and 9)
    GJ_L = self.G * self.J / L
    k_local[3, 3] = GJ_L
    k_local[3, 9] = -GJ_L
    k_local[9, 3] = -GJ_L
    k_local[9, 9] = GJ_L
    
    # Bending in local XY plane (about Z axis) - DOF 1, 5, 7, 11
    EIz_L3 = self.E * self.Iz / (L**3)
    k_local[1, 1] = 12 * EIz_L3
    k_local[1, 5] = 6 * EIz_L3 * L
    k_local[1, 7] = -12 * EIz_L3
    k_local[1, 11] = 6 * EIz_L3 * L
    
    k_local[5, 1] = 6 * EIz_L3 * L
    k_local[5, 5] = 4 * EIz_L3 * L**2
    k_local[5, 7] = -6 * EIz_L3 * L
    k_local[5, 11] = 2 * EIz_L3 * L**2
    
    # ... (similar terms for other DOF)
    
    return k_local
\end{lstlisting}

\subsection{Συναρμολόγηση Ολικού Συστήματος}

Ο ολικός πίνακας δυσκαμψίας συναρμολογείται μετασχηματίζοντας τον τοπικό πίνακα δυσκαμψίας κάθε στοιχείου σε ολικές συντεταγμένες και προσθέτοντας τις συνεισφορές:

\begin{lstlisting}[caption={Συναρμολόγηση Ολικού Πίνακα Δυσκαμψίας (Μερική)}]
def assemble_stiffness_matrix(self):
    """Assemble global stiffness matrix for beam elements"""
    print("\nASSEMBLING GLOBAL STIFFNESS MATRIX (BEAM ELEMENTS)")
    print("="*80)
    
    num_nodes = len(self.nodes)
    n_dof = self.dof_per_node * num_nodes  # 6 DOF per node
    
    self.K_global = np.zeros((n_dof, n_dof))
    
    # Loop through all elements
    for idx in range(len(self.elements)):
        elem = self.elements.iloc[idx]
        
        node1_num = int(elem['Node1'])
        node2_num = int(elem['Node2'])
        
        # Get node coordinates
        node1 = self.nodes[self.nodes['Node Number'] == node1_num].iloc[0]
        node2 = self.nodes[self.nodes['Node Number'] == node2_num].iloc[0]
        
        x1, y1, z1 = node1['X'], node1['Y'], node1['Z']
        x2, y2, z2 = node2['X'], node2['Y'], node2['Z']
\end{lstlisting}

\subsection{Επίλυση Συστήματος}

Μετά την εφαρμογή των οριακών συνθηκών, το σύστημα επιλύεται χρησιμοποιώντας τον επιλύτη γραμμικής άλγεβρας του NumPy:

\begin{lstlisting}[caption={Επίλυση Συστήματος}]
def solve(self):
    """Solve the FEM system"""
    print("\nSOLVING FEM SYSTEM")
    print("="*80)
    
    # Identify fixed DOFs
    fixed_dofs = []
    for node_id, bc in self.boundary_conditions.items():
        node_idx = node_id - 1  # Convert to 0-based
        if bc['Ux'] == 1:
            fixed_dofs.append(node_idx * self.dof_per_node + 0)
        if bc['Uy'] == 1:
            fixed_dofs.append(node_idx * self.dof_per_node + 1)
        if bc['Uz'] == 1:
            fixed_dofs.append(node_idx * self.dof_per_node + 2)
        if bc['Rx'] == 1:
            fixed_dofs.append(node_idx * self.dof_per_node + 3)
        if bc['Ry'] == 1:
            fixed_dofs.append(node_idx * self.dof_per_node + 4)
        if bc['Rz'] == 1:
            fixed_dofs.append(node_idx * self.dof_per_node + 5)
    
    fixed_dofs = np.array(fixed_dofs, dtype=int)
    all_dofs = np.arange(len(self.F_global))
    free_dofs = np.setdiff1d(all_dofs, fixed_dofs)
    
    # Reduce system to free DOFs
    K_reduced = self.K_global[np.ix_(free_dofs, free_dofs)]
    F_reduced = self.F_global[free_dofs]
    
    # Check conditioning
    cond = np.linalg.cond(K_reduced)
    print(f"  Condition number: {cond:.2e}")
    
    if cond > 1e10:
        print("  WARNING: Stiffness matrix is ill-conditioned")
    
    # Solve for free DOFs
    u_reduced = np.linalg.solve(K_reduced, F_reduced)
    
    # Expand to full displacement vector
    self.u_global = np.zeros(len(self.F_global))
    self.u_global[free_dofs] = u_reduced
    
    print(f"✓ System solved successfully")
\end{lstlisting}

\section{Μονάδα Μετεπεξεργαστή}

Ο μετεπεξεργαστής (\texttt{postprocessor.py}) παρέχει δυνατότητες διαδραστικής οπτικοποίησης.

\subsection{Δομή Κλάσης}

\begin{lstlisting}[caption={Αρχικοποίηση Κλάσης FEMPostProcessor}]
class FEMPostProcessor:
    """Handles visualization of FEM results"""
    
    def __init__(self):
        self.nodes = pd.DataFrame()
        self.elements = pd.DataFrame()
        self.displacements = pd.DataFrame()
        self.reactions = pd.DataFrame()
        self.element_forces = pd.DataFrame()
        self.summary = {}
\end{lstlisting}

\subsection{Οπτικοποίηση Παραμορφωμένης Κατασκευής}

Ο μετεπεξεργαστής δημιουργεί διαδραστικές τρισδιάστατες οπτικοποιήσεις με ρυθμιζόμενη μεγέθυνση:

\begin{lstlisting}[caption={Γράφημα Παραμορφωμένης Κατασκευής με Slider (Μερικό)}]
def plot_deformed_structure(self, magnification=1.0, show_plot=True):
    """Plot undeformed and deformed structure with interactive magnification slider"""
    print("\nCREATING DEFORMATION PLOT WITH INTERACTIVE SLIDER")
    print("="*80)
    
    fig = go.Figure()
    
    # Undeformed structure (always shown)
    edge_x_undef, edge_y_undef, edge_z_undef = [], [], []
    for idx in range(len(self.elements)):
        node1_num = int(self.elements.iloc[idx]['Node1'])
        node2_num = int(self.elements.iloc[idx]['Node2'])
        
        node1 = self.nodes[self.nodes['Node Number'] == node1_num].iloc[0]
        node2 = self.nodes[self.nodes['Node Number'] == node2_num].iloc[0]
        
        edge_x_undef.extend([node1['X'], node2['X'], None])
        edge_y_undef.extend([node1['Y'], node2['Y'], None])
        edge_z_undef.extend([node1['Z'], node2['Z'], None])
    
    fig.add_trace(go.Scatter3d(
        x=edge_x_undef, y=edge_y_undef, z=edge_z_undef,
        mode='lines',
        line=dict(color='cyan', width=4),
        name='Undeformed',
        visible=True,
        opacity=0.9
    ))
    
    # Create deformed structures for different magnifications
    magnifications = [1, 10, 50, 100, 500, 1000, 2000, 5000]
    
    for mag in magnifications:
        # Deformed structure
        deformed_nodes = self.nodes.copy()
        for idx in range(len(self.nodes)):
            node_num = int(self.nodes.iloc[idx]['Node Number'])
            disp = self.displacements[self.displacements['Node Number'] == node_num].iloc[0]
            
            deformed_nodes.loc[idx, 'X'] += disp['Ux'] * mag
            deformed_nodes.loc[idx, 'Y'] += disp['Uy'] * mag
            deformed_nodes.loc[idx, 'Z'] += disp['Uz'] * mag
\end{lstlisting}

\subsection{Οπτικοποίηση Τάσεων και Παραμορφώσεων}

Ο μετεπεξεργαστής δημιουργεί επίσης οπτικοποίηση των κατανομών τάσεων και παραμορφώσεων:

\begin{lstlisting}[caption={Γράφημα Κατανομής Τάσεων (Μερικό)}]
def plot_stress_distribution(self, show_plot=True):
    """Plot stress distribution on elements"""
    print("\nCREATING STRESS DISTRIBUTION PLOT")
    print("="*80)
    
    fig = go.Figure()
    
    # Get stress values
    stresses = self.element_forces['Stress'].values
    max_abs_stress = np.max(np.abs(stresses))
    
    # Color mapping
    edge_x, edge_y, edge_z, edge_colors = [], [], [], []
    
    for idx in range(len(self.elements)):
        node1_num = int(self.elements.iloc[idx]['Node1'])
        node2_num = int(self.elements.iloc[idx]['Node2'])
        
        node1 = self.nodes[self.nodes['Node Number'] == node1_num].iloc[0]
        node2 = self.nodes[self.nodes['Node Number'] == node2_num].iloc[0]
        
        stress = stresses[idx]
        
        # Create line segment for each element
        edge_x.extend([node1['X'], node2['X'], None])
        edge_y.extend([node1['Y'], node2['Y'], None])
        edge_z.extend([node1['Z'], node2['Z'], None])
        edge_colors.extend([stress, stress, None])
\end{lstlisting}

\section{Κύριο Σενάριο Εκτέλεσης}

Το σενάριο \texttt{FEM\_Main.py} συντονίζει ολόκληρη τη διαδικασία ανάλυσης:

\begin{lstlisting}[caption={Πλήρης Διαδικασία Ανάλυσης FEM}]
def main():
    """Execute the complete FEM analysis pipeline"""
    
    print("\n" + "="*80)
    print("FEM ANALYSIS - COMPLETE PIPELINE")
    print("="*80 + "\n")
    
    # Step 1: Pre-processor
    print("STEP 1: Running Pre-Processor...")
    print("-"*80)
    import preprocessor
    preprocessor.main()
    
    print("\n" + "="*80 + "\n")
    
    # Step 2: Solver
    print("STEP 2: Running Solver...")
    print("-"*80)
    import solver
    solver.main()
    
    print("\n" + "="*80 + "\n")
    
    # Step 3: Post-processor
    print("STEP 3: Running Post-Processor...")
    print("-"*80)
    import postprocessor
    postprocessor.main()
    
    print("\n" + "="*80)
    print("FEM ANALYSIS COMPLETED SUCCESSFULLY")
    print("="*80 + "\n")
    
    print("Results available in:")
    print("  - data/structure.dat (input data)")
    print("  - data/structure.res (results)")
    print("  - plots/ (interactive visualizations)")
    print("  - output_html/ (data tables)")
    print("\nOpen the HTML files in plots/ folder to view interactive results.")
\end{lstlisting}

\section{Μορφές Αρχείων Δεδομένων}

\subsection{Μορφή Αρχείου Εισόδου (.dat)}

Το αρχείο ορισμού κατασκευής έχει την ακόλουθη μορφή:

\begin{verbatim}
PARAMETERS
E A Iz Iy J G diameter

MATERIALS
STEEL E nu

NODES
node_id x y z

ELEMENTS
elem_id node1 node2 cross_section E nu

BOUNDARY_CONDITIONS
node_id Ux Uy Uz Rx Ry Rz

LOADS
node_id Fx Fy Fz Mx My Mz

END
\end{verbatim}

\subsection{Μορφή Αρχείου Αποτελεσμάτων (.res)}

Το αρχείο αποτελεσμάτων περιέχει:

\begin{verbatim}
DISPLACEMENTS
node_id Ux Uy Uz Rx Ry Rz |U| |R|

REACTIONS
node_id Rx Ry Rz Mx My Mz

ELEMENT_FORCES
elem_id Axial Shear_Y Shear_Z Torsion Moment_Y Moment_Z Stress Strain

SUMMARY
max_displacement value
max_stress value
max_strain value

END
\end{verbatim}

\section{Παράδειγμα Χρήσης}

Για να χρησιμοποιήσετε το λογισμικό για μια προσαρμοσμένη κατασκευή:

\begin{enumerate}
    \item \textbf{Ορίστε Γεωμετρία}: Επεξεργαστείτε τις μεθόδους \texttt{\_define\_geometry\_parameters()} και \texttt{\_add\_initial\_nodes()} στο \texttt{preprocessor.py}
    
    \item \textbf{Ορίστε Οριακές Συνθήκες}: Προσθέστε σταθερούς κόμβους χρησιμοποιώντας \texttt{add\_boundary\_condition()}
    
    \item \textbf{Εφαρμόστε Φορτία}: Ορίστε φορτία χρησιμοποιώντας \texttt{add\_load()}
    
    \item \textbf{Εκτελέστε Ανάλυση}: Εκτελέστε \texttt{python FEM\_Main.py}
    
    \item \textbf{Δείτε Αποτελέσματα}: Ανοίξτε τα δημιουργημένα αρχεία HTML στον κατάλογο \texttt{plots/}
\end{enumerate}

\section{Επικύρωση}

Ο κώδικας έχει επικυρωθεί έναντι αναλυτικών λύσεων για απλές κατασκευές. Για παράδειγμα, μια τετραεδρική κατασκευή πλαισίου αναλύθηκε και τα αποτελέσματα αντιστοιχούσαν στην αναλυτική λύση που τεκμηριώνεται στο αρχείο \texttt{README.md}.

\subsection{Δοκιμαστική Περίπτωση: Απλή Τετραεδρική Κατασκευή}

\textbf{Γεωμετρία:}
\begin{itemize}
    \item 4 κόμβοι που σχηματίζουν τετράεδρο
    \item 6 στοιχεία ράβδων
    \item 3 σταθεροί κόμβοι, 1 φορτισμένος κόμβος
\end{itemize}

\textbf{Ιδιότητες Υλικού:}
\begin{itemize}
    \item $E = 210$ GPa (Χάλυβας)
    \item $A = 4$ cm$^2$
\end{itemize}

\textbf{Φόρτιση:}
\begin{itemize}
    \item $F_x = 1000$ N σε ελεύθερο κόμβο
\end{itemize}

Οι υπολογισμένες μετατοπίσεις και τάσεις αντιστοιχούν στην αναλυτική λύση εντός αριθμητικής ακρίβειας.

\section{Συμπεράσματα}

Αυτή η υλοποίηση λογισμικού FEM παρέχει μια πλήρη λύση για την ανάλυση τρισδιάστατων κατασκευών δοκών. Τα κύρια επιτεύγματα περιλαμβάνουν:

\begin{itemize}
    \item Ισχυρή διατύπωση τρισδιάστατου στοιχείου δοκού με 6 DOF ανά κόμβο
    \item Αρθρωτή αρχιτεκτονική που διευκολύνει τη συντήρηση και επέκταση κώδικα
    \item Αυτόματη δημιουργία στοιχείων για σύνθετες γεωμετρίες
    \item Ολοκληρωμένες δυνατότητες οπτικοποίησης
    \item Επικυρωμένο έναντι αναλυτικών λύσεων
\end{itemize}

\subsection{Μελλοντικές Βελτιώσεις}

Πιθανές βελτιώσεις για μελλοντικές εκδόσεις:
\begin{itemize}
    \item Υποστήριξη για μη κυκλικές διατομές
    \item Μη γραμμικότητα υλικού
    \item Δυνατότητες δυναμικής ανάλυσης
    \item Αλγόριθμοι βελτίωσης πλέγματος
    \item Διεπαφή GUI για ορισμό γεωμετρίας
\end{itemize}

\section{Αναφορές}

\begin{enumerate}
    \item Bathe, K. J. (1996). \textit{Finite Element Procedures}. Prentice Hall.
    \item Zienkiewicz, O. C., \& Taylor, R. L. (2000). \textit{The Finite Element Method}. Butterworth-Heinemann.
    \item Cook, R. D., Malkus, D. S., Plesha, M. E., \& Witt, R. J. (2001). \textit{Concepts and Applications of Finite Element Analysis}. Wiley.
\end{enumerate}

\appendix

\section{Πλήρεις Λίστες Κώδικα}

Για τις πλήρεις λεπτομέρειες υλοποίησης, ανατρέξτε στα ακόλουθα αρχεία:
\begin{itemize}
    \item \texttt{preprocessor.py} (752 γραμμές)
    \item \texttt{solver.py} (694 γραμμές)
    \item \texttt{postprocessor.py} (732 γραμμές)
    \item \texttt{geometry\_utils.py} (29 γραμμές)
    \item \texttt{FEM\_Main.py} (56 γραμμές)
\end{itemize}

\section{Εξαρτήσεις Λογισμικού}

Η υλοποίηση απαιτεί τα ακόλουθα πακέτα Python:
\begin{itemize}
    \item \texttt{numpy}: Αριθμητικοί υπολογισμοί
    \item \texttt{pandas}: Διαχείριση δομών δεδομένων
    \item \texttt{plotly}: Διαδραστική οπτικοποίηση
\end{itemize}

Εγκαταστήστε τις εξαρτήσεις χρησιμοποιώντας:
\begin{verbatim}
pip install numpy pandas plotly
\end{verbatim}

\end{document}
